\chapter{Stable lattices and reduction}
\label{chap:library:reduction}

A stable lattice in an ordinary representation gives a modular
representation by reduction modulo the maximal ideal. Different lattices
may have nonisomorphic reductions, but the multiplicities of their simple
composition factors agree. The decomposition homomorphism assigns this
common Grothendieck class to the ordinary representation.

The construction below identifies tensor reduction with reduction modulo
the maximal ideal. Lattice independence and the decomposition homomorphism
then follow from the ordinary and Brauer character identity in
Hypothesis~\ref{lib:reduction:character-hypothesis}.

Throughout the chapter, $G$ is a finite group and $\ell$ is prime.
The coefficient fields for ordinary and modular representations are
$K$ and $k$, respectively. The integral coefficient ring is denoted by
$\mathcal O$. Each section states the coefficient hypotheses it uses.

\section{Integral coefficients and the residue map}
\label{lib:reduction:modular-system}

\begin{definition}
A complete $\ell$-modular system $(K,\mathcal O,k)$ consists of a complete
discrete valuation ring $\mathcal O$, its fraction field $K$ of
characteristic zero, a field $k$ of characteristic $\ell$, and a specified
isomorphism
\[
 \mathcal O/\mathfrak m\simeq k,
\]
where $\mathfrak m$ is the maximal ideal of $\mathcal O$. Completeness is
with respect to the $\mathfrak m$-adic topology. The residue map
$\pi:\mathcal O\to k$ is the composite of the quotient map with this
isomorphism.
\end{definition}

The map $\pi$ is onto and has kernel $\mathfrak m$. Its choice specifies
the $\mathcal O$-algebra structure on $k$ used in tensor reduction.
Splitting hypotheses for $G$ and a correspondence between roots of unity
in $k$ and $K$ are additional conditions. Algebraic closure of the
fraction field $K$ is not required.

\begin{implementation}
The definition in \module{ModularRep.ModularSystem} includes the prime,
the complete discrete valuation ring, its fraction field, the two
characteristics, and the residue field isomorphism \lean{residueEquiv}.
The definitions \lean{ModularSystem.residue} and
\lean{ModularSystem.residueAlgebra} give $\pi$ and the induced
$\mathcal O$-algebra structure on $k$.
The theorems \lean{residue_surjective} and \lean{ker_residue} prove the
two properties of $\pi$ used below. All reductions associated with this
modular system use \lean{residueAlgebra}.
\end{implementation}

\section{Constructing a stable lattice}
\label{lib:reduction:stable-lattice}

\begin{samepage}
For this section it is enough that $\mathcal O$ is a commutative ring,
$K$ is a field with an $\mathcal O$-algebra structure, and $V$ is a finite
dimensional $K$-space, with its compatible $\mathcal O$-module structure.
Let $\rho:G\to\GL_K(V)$ be a representation.

\end{samepage}

\begin{definition}
An $\mathcal O$-lattice in $V$ is a finitely generated
$\mathcal O$-submodule $L\leq V$ whose $K$-linear span is $V$.
It is stable if $\rho(g)L\subseteq L$ for every $g\in G$.
\end{definition}

When $\mathcal O$ is a discrete valuation ring with fraction field $K$,
this agrees with the usual notion of a full lattice. Only finite
generation and full span are needed for the construction below. For the
usual construction over a discrete valuation ring, see, for example,
\cite[\S15.2]{Serre1977}.

\begin{proposition}[Stable closure of a lattice]
\label{lib:reduction:existence}
If $L_0$ is an $\mathcal O$-lattice in $V$, then
\[
 L=\sum_{g\in G}\rho(g)L_0
\]
is a stable $\mathcal O$-lattice containing $L_0$. In particular every
finite dimensional $K$-representation of a finite group has a stable
$\mathcal O$-lattice.
\end{proposition}
\begin{proof}
Each image $\rho(g)L_0$ is finitely generated over $\mathcal O$.
There are only finitely many such submodules, so their sum is finitely
generated. The term $g=1$ shows $L_0\subseteq L$, which implies that the
$K$-span of $L$ is all of $V$. For $h\in G$,
$\rho(h)\rho(g)L_0=\rho(hg)L_0$ is another summand. Thus $\rho(h)$
preserves $L$.

For existence, choose a finite $K$-basis $b_1,\ldots,b_n$ of $V$ and
take $L_0=\sum_i\mathcal O b_i$. This submodule is finitely generated
and spans $V$ over $K$, so the construction applies.
\end{proof}

\begin{implementation}
\module{ModularRep.StableLattice} defines \lean{StableLattice} as a stable
subrepresentation satisfying Mathlib's \lean{IsLattice} predicate. The orbit sum is
\lean{StableLattice.orbitSpan}, represented by the supremum of the image
submodules. Its containment, stability, finite generation, and full span
are proved in \lean{le_orbitSpan}, \lean{orbitSpan_stable},
\lean{orbitSpan_fg}, and \lean{orbitSpan_isLattice}. The resulting lattice
is \lean{stableClosure}, and the basis construction is \lean{ofBasis}.
The construction \lean{ofBasis} also applies when $K$ is a commutative
ring and a finite basis is given. It requires neither a discrete valuation
nor completeness.
\end{implementation}

The integral representation on $L$ is obtained by restricting $\rho(g)$
to $L$ for each $g$. Once an $\mathcal O$-algebra structure on $k$ is
specified, define its reduction by
\[
 \overline L=k\otimes_{\mathcal O}L,
 \qquad
 g(a\otimes x)=a\otimes\rho(g)x.
\]
If $x_1,\ldots,x_t$ generate $L$ over $\mathcal O$, the tensors
$1\otimes x_1,\ldots,1\otimes x_t$ generate $\overline L$ over $k$.
Indeed expansion in the generators and the balancing relation move all
$\mathcal O$-coefficients to the first tensor factor. Hence the reduction
is finite dimensional whenever $k$ is a field.

\begin{implementation}
The declarations \lean{StableLattice.integralRepresentation} and
\lean{StableLattice.reduction} construct these actions.
\lean{reduction_apply_tmul} gives the action on pure tensors, and
\lean{reduction_moduleFinite} proves finite generation by the tensors of
an arbitrary finite generating set. In
\module{ModularRep.StableLatticeKZero}, \lean{chosenStableLattice}
uses a finite basis of an \lean{FDRep} object, and
\lean{stableLatticeReduction} is its finite dimensional reduction.
\end{implementation}

\section{Tensor reduction as an equivariant quotient}
\label{lib:reduction:quotient}

The quotient description of reduction is valid under weaker hypotheses
than those of a modular system. Let $\pi:\mathcal O\to k$ be a
surjective homomorphism of commutative rings, and write $I=\ker\pi$.
For an $\mathcal O$-module $M$, the notation $IM$ means the submodule
generated by the products $rx$ with $r\in I$ and $x\in M$.
The quotient $M/IM$ is a $k$-module through
$\mathcal O/I\simeq k$. Write $\overline x=x+IM$ for the class of $x\in M$.

\begin{proposition}[Reduction modulo an ideal]
\label{lib:reduction:tensor-quotient-theorem}
There is a $k$-linear isomorphism
\[
 \Phi:k\otimes_{\mathcal O}M\xrightarrow{\sim}M/IM,
 \qquad a\otimes x\longmapsto a\overline x,
 \qquad \Phi^{-1}(\overline x)=1\otimes x.
\]
If $M$ carries an $\mathcal O$-linear action of $G$, then $IM$ is stable,
the action descends by $g\overline x=\overline{gx}$, and $\Phi$ intertwines
the tensor and quotient actions.
\end{proposition}
\begin{proof}
Surjectivity gives the isomorphism $\mathcal O/I\simeq k$.
The standard tensor quotient isomorphism identifies
$(\mathcal O/I)\otimes_{\mathcal O}M$ with $M/IM$. Using
$\mathcal O/I\simeq k$ in the first tensor factor gives $\Phi$ with the stated
formula. The quotient relation makes $\overline x\mapsto1\otimes x$
well defined: if $r\in I$, then
$1\otimes rx=\pi(r)\otimes x=0$.

Every group operator is $\mathcal O$-linear, so it sends $rx$ to $r(gx)$
and preserves $IM$. This defines the quotient action directly. For a pure
tensor,
\[
 \Phi\bigl(g(a\otimes x)\bigr)=a\overline{gx}
   =g(a\overline x)=g\Phi(a\otimes x).
\]
Additivity extends the identity to the whole tensor product.
\end{proof}

\begin{implementation}
\module{ModularRep.ReductionModulo} defines $IM$ as
\lean{ReductionModulo.kernelSMul}. Its \lean{tensorQuotientEquiv}
combines Mathlib's \lean{TensorProduct.quotTensorEquivQuotSMul} with
its first isomorphism theorem for algebras. The local construction
transports the quotient scalar action and uses surjectivity to obtain
the $k$-linear equivalence. The versions for an explicit ring map and a
specified ideal are \lean{tensorQuotientEquivOfSurjectiveRingHom} and
\lean{tensorQuotientEquivOfKernelEq}.

\module{ModularRep.ReductionModuloEquivariance} first constructs
\lean{ReductionModulo.quotientRepresentation} by descending the integral
action. The intertwining identity is
\lean{tensorQuotientEquiv_isIntertwining}, and the resulting equivalence
of representations is \lean{tensorQuotientRepresentationEquiv}.
\end{implementation}

For a modular system, $I=\mathfrak m$ and the proposition gives
$\overline L\simeq L/\mathfrak mL$ as $kG$-modules. This is
\lean{stableLatticeReductionQuotientEquiv}. Its underlying linear
equivalence agrees with
\lean{stableLatticeReductionCarrierQuotientEquiv}.
For a fixed stable lattice, this isomorphism does not require $G$ to be finite.

\section{Character compatibility and lattice independence}
\label{lib:reduction:compatibility}

Fix a complete $\ell$-modular system. For the character construction in
this section, additionally assume that $k$ is algebraically closed and
choose a finite root correspondence $\iota$ for $G$. The field $K$ has
characteristic zero and is the fraction field of $\mathcal O$. Write
\[
 \operatorname{br}_{\iota}:K_0(kG)\longrightarrow
     \operatorname{CF}_{\ell'}(G,K),
 \qquad
 c:K_0(KG)\longrightarrow\operatorname{CF}_{\ell'}(G,K)
\]
for the Brauer character homomorphism and the homomorphism taking an
ordinary representation class to its trace character restricted to
$\ell$-regular elements. The first is injective by
Proposition~\ref{lib:kzero:brauer-injective}. The second is defined by
ordinary trace additivity and the universal property of the Grothendieck group.

\begin{hypothesis}[Reduction compatibility]
\label{lib:reduction:character-hypothesis}
Write $G_{\ell'}$ for the set of $\ell$-regular elements. For every finite
dimensional $K$-representation $V$, with ordinary trace character $\chi_V$,
and every stable $\mathcal O$-lattice $L$ in $V$, the fixed root
correspondence satisfies
\[
 \varphi_{k\otimes_{\mathcal O}L}
      =\chi_V\big|_{G_{\ell'}}.
\]
\end{hypothesis}

The reduction identity with compatible roots is standard
\cite[Corollary~2.9]{Navarro1998}.
Here the residue map and the root correspondence are specified
separately. Their compatibility is required: the identity is assumed
for every stable lattice with these same choices.

\begin{implementation}
This hypothesis is \lean{StableReductionBrauerCharacterCompatibility}.
It is an additional condition on \lean{ModularSystem} and
\lean{PrimeRegularRootEmbedding}.
\end{implementation}

Lattice independence and the decomposition homomorphism are treated in
\cite[\S15.2, Theorem~32 and the following discussion]{Serre1977}. That account proves lattice independence by
comparing lattices directly. It also notes that completeness is unnecessary
for that argument. Under the coefficient hypotheses and character identity
stated above, lattice independence also follows from injectivity of the
Brauer character homomorphism, as in the proof below.

\begin{theorem}[Decomposition homomorphism]
\label{lib:reduction:decomposition-theorem}
Under Hypothesis~\ref{lib:reduction:character-hypothesis}, there is a unique
homomorphism of abelian groups
\[
 d:K_0(KG)\longrightarrow K_0(kG)
 \quad\text{such that}\quad
 \operatorname{br}_{\iota}\circ d=c.
\]
For every stable lattice $L$ in $V$, it satisfies
$d([V])=[k\otimes_{\mathcal O}L]$. Consequently two stable lattices in
the same ordinary representation have equal modular Grothendieck classes.
\end{theorem}
\begin{proof}
Choose one stable lattice for every ordinary representation, using
Proposition~\ref{lib:reduction:existence}, and assign its reduction class
to the representation. If $L_1,L_2$ lie in the same $V$, compatibility gives
\[
 \operatorname{br}_{\iota}([\overline L_1])
  =c([V])
  =\operatorname{br}_{\iota}([\overline L_2]).
\]
Injectivity implies $[\overline L_1]=[\overline L_2]$. The same argument
using invariance of $c$ under isomorphism compares choices for isomorphic
ordinary representations.

For a short exact sequence with terms $U,V,W$, let $a_U,a_V,a_W$ be the
chosen reduction classes. Ordinary trace additivity gives
\[
 \operatorname{br}_{\iota}(a_V)
   =c([V])=c([U])+c([W])
   =\operatorname{br}_{\iota}(a_U+a_W).
\]
Injectivity gives $a_V=a_U+a_W$.
The universal property therefore defines $d$. On every object class,
compatibility gives $\operatorname{br}_{\iota}d=c$, so the two
homomorphisms agree everywhere. Comparing $d([V])$ with any other stable
lattice reduction proves the asserted formula for every $L$.
Finally, if $d'$ has the same character identity, then
$\operatorname{br}_{\iota}(d'(x))=\operatorname{br}_{\iota}(d(x))$
for every $x$. Injectivity gives $d'=d$.
\end{proof}

\begin{implementation}
\module{ModularRep.ExactCharacterBridge} proves the same assertion for
homomorphisms from the two Grothendieck groups to an abelian group $A$,
provided the modular homomorphism is injective and the character identity
holds for every stable lattice. The resulting definitions are
\lean{ExactCharacterBridge.toExactDecompositionMapData} and
\lean{ExactCharacterBridge.decompositionMap}.

\module{ModularRep.BrauerDecompositionMap} specialises this argument to
ordinary and Brauer characters. Its
\lean{exactCharacterBridgeOfStableReduction} uses
\lean{brauerCharacterKZeroHom_injective} for injectivity on $K_0(kG)$.
The resulting homomorphism is \lean{decompositionMapOfStableReduction},
its character identity is
\lean{decompositionMapOfStableReduction_characterSquare}, and its
uniqueness is \lean{existsUnique_decompositionMapOfStableReduction}.
The declarations
\lean{decompositionMap_classOf_eq_stableLatticeReduction} and
\lean{stableLatticeReductionClass_eq} give the two lattice statements.
\end{implementation}

By Theorem~\ref{lib:kzero:basis-theorem}, equality of these
classes is equivalent to equality of their composition multiplicities.

\begin{example}[Decomposition numbers]
\label{lib:reduction:coefficient-example}
For a simple $kG$-module $S$, consider the coefficient of $[S]$ in
$d([V])$ with respect to the basis of simple module classes. If $L$ is any
stable lattice in $V$, this coefficient is $m_S(\overline L)$ and is a nonnegative integer.
For a short exact sequence $0\to U\to V\to W\to0$, it is the sum of
the coefficients for $U$ and $W$. For irreducible $V$, these coefficients
are the decomposition numbers.
This definition makes sense over the coefficient fields fixed here. To
identify them with the usual character decomposition numbers for $G$, one
also chooses $K$ to be a splitting field for $G$.
\end{example}

\section{Twists and tensor products}
\label{lib:reduction:naturality}

For $\alpha\in\Aut(G)$, define the twist of $V$ by precomposition:
$\rho^{\alpha}(g)=\rho(\alpha(g))$. The induced operation on exact
$K_0$ is denoted by $T_{\alpha}$. It comes from an exact autoequivalence,
and the same notation is used over $K$ and $k$.

\begin{proposition}[Compatibility with automorphisms]
\label{lib:reduction:twist-theorem}
Under the hypotheses of Theorem~\ref{lib:reduction:decomposition-theorem},
\[
 d(T_{\alpha}x)=T_{\alpha}d(x)
 \qquad (\alpha\in\Aut(G),\ x\in K_0(KG)).
\]
\end{proposition}
\begin{proof}
Both character homomorphisms commute with precomposition by $\alpha$.
Applying $\operatorname{br}_{\iota}$ to the two sides therefore gives
the same function, namely $c(x)\circ\alpha$ on $G_{\ell'}$.
Injectivity proves the equality in $K_0(kG)$.
\end{proof}

\begin{implementation}
\module{ModularRep.KZeroTwist} defines \lean{FDRep.twistEquivalence}
by applying Mathlib's \lean{Action.resEquiv} to finite dimensional
representations. It constructs the induced map \lean{twistKZero} and proves
\lean{ordinaryCharacterKZero_twist} and
\lean{brauerCharacterKZeroHom_twist}. The consequence above is
\lean{decompositionMapOfStableReduction_twist}.
The convention is a right action: with automorphism multiplication given
by composition, \lean{twistKZero_comp} states
$T_{\alpha}\circ T_{\beta}=T_{\beta\alpha}$.
\end{implementation}

For a linear character $\lambda:G\to k^{\times}$, write $L_{\lambda}V$
for the representation on the same vector space with action
$g\cdot v=\lambda(g)\rho(g)v$. This is tensoring by the corresponding
one dimensional representation. Tensoring by $\lambda^{-1}$ is its
inverse, so it is an exact autoequivalence. Use $L_{\lambda}$ also for
the induced operation on $K_0$.

Twisting and tensoring satisfy
\[
 T_\alpha L_\lambda=L_{\lambda\circ\alpha}T_\alpha.
\]
The two composite exact functors are naturally isomorphic: the identity
on the representation space intertwines their equal actions. This
argument applies over either coefficient field and requires neither a
modular system nor a formula for Brauer characters of tensor products.
The natural isomorphism is
\lean{ModularRep.FDRep.linearCharacterTwistTwistIso}, and its consequence
on $K_0$ is \lean{ModularRep.twistKZero_linearCharacterTwistKZero}, both
in \module{ModularRep.KZeroLinearCharacterSemidirect}.

For a modular linear character $\lambda$, write $\widehat\lambda(g)$
for the lift of $\lambda(g)$ through $\iota$, for $g\in G_{\ell'}$.
Let $\lambda_K:G\to K^{\times}$ and $\lambda_k:G\to k^{\times}$
be linear characters, and assume that
$\lambda_K\big|_{G_{\ell'}}=\widehat{\lambda_k}$.
By Theorem~\ref{lib:brauer:tensor-source},
\[
 \varphi_{L_{\lambda}V}=\widehat\lambda\,\varphi_V
 \qquad(V\in\operatorname{FDRep}_k(G),\ \lambda:G\to k^{\times})
\]
on $G_{\ell'}$, with pointwise multiplication on the right.
For the general tensor product formula, see
\cite[proof of Theorem~2.23]{Navarro1998}.

\begin{proposition}[Compatibility with linear characters]
\label{lib:reduction:tensor-theorem}
If $d$ is the homomorphism of
Theorem~\ref{lib:reduction:decomposition-theorem} and $\lambda_K$ and
$\lambda_k$ satisfy the compatibility condition above, then, for every
$x\in K_0(KG)$,
\[
 d(L_{\lambda_K}x)=L_{\lambda_k}d(x).
\]
Consequently the combined operation also satisfies
$d(L_{\lambda_K}T_{\alpha}x)=L_{\lambda_k}T_{\alpha}d(x)$.
\end{proposition}
\begin{proof}
Ordinary trace on $L_{\lambda_K}V$ is the product of $\lambda_K$ with
the trace on $V$, since scalar multiplication by $\lambda_K(g)$ scales
the trace by the same factor. Theorem~\ref{lib:brauer:tensor-source}
gives the analogous statement over $k$. Compatibility of the two linear characters
therefore identifies the images under $\operatorname{br}_{\iota}$ of
$d(L_{\lambda_K}x)$ and $L_{\lambda_k}d(x)$. Injectivity gives the first
identity. Combining it with Proposition~\ref{lib:reduction:twist-theorem}
gives the second.
\end{proof}

\begin{implementation}
\module{ModularRep.KZeroLinearCharacterTensor} constructs the exact
autoequivalence and \lean{linearCharacterTwistKZero}, and proves the
ordinary trace formula. In
\module{ModularRep.KZeroLinearCharacterDecomposition},
\lean{LinearCharacterReductionCompatible} states the equality
$\lambda_K\big|_{G_{\ell'}}=\widehat{\lambda_k}$. The conclusions are
\lean{decompositionMapOfStableReduction_linearCharacterTwist} and
\lean{decompositionMapOfStableReduction_linearCharacterTwist_twist}.
Both the compatibility of the linear characters and
\lean{BrauerLinearTensorProductFormula} are hypotheses of these results.
\end{implementation}

\section{Specifying a decomposition homomorphism}
\label{lib:reduction:interfaces}

An assignment $[V]\mapsto a_V\in K_0(kG)$ that is additive on short
exact sequences induces a homomorphism $K_0(KG)\to K_0(kG)$. If
$a_V=[\overline L]$ for every stable lattice $L$ in $V$, this is the
decomposition homomorphism.
Theorem~\ref{lib:reduction:decomposition-theorem} proves these properties
from character compatibility.

\begin{implementation}
\lean{ExactDecompositionMapData} describes an assignment on ordinary
isomorphism classes that is additive on short exact sequences. Its
\lean{toAddMonoidHom} is the homomorphism given by the universal property.
The condition \lean{RealisesStableLatticeReduction} says that every
stable lattice has the assigned reduction class. For
\lean{decompositionMapOfStableReduction}, this condition is proved by
\lean{exactDecompositionMapDataOfStableReduction_realises}.

\lean{DecompositionMapInterface} describes a homomorphism between
specified ordinary and modular class groups that agrees with reduction
for every stable lattice. Equality of the classes from
two lattices then follows by expressing both as the image of the same
ordinary class.
\end{implementation}
